\documentclass[11pt,a4paper]{article}

% === Packages ===
\usepackage[utf8]{inputenc}
\usepackage[T1]{fontenc}
\usepackage{amsmath,amssymb,amsthm,mathtools}
\usepackage{enumitem}
\usepackage{hyperref}
\usepackage{cleveref}
\usepackage{tikz}
\usetikzlibrary{positioning,arrows.meta,calc}
\usepackage{booktabs}
\usepackage{geometry}
\usepackage{xcolor}
\usepackage{tcolorbox}
\usepackage{fancyvrb}
\usepackage{listings}

\geometry{margin=1in}
\hypersetup{colorlinks=true,linkcolor=blue!70!black,citecolor=green!50!black,urlcolor=purple!70!black}

% === Theorem environments ===
\theoremstyle{plain}
\newtheorem{theorem}{Theorem}[section]
\newtheorem{proposition}[theorem]{Proposition}
\newtheorem{lemma}[theorem]{Lemma}
\newtheorem{corollary}[theorem]{Corollary}
\newtheorem{conjecture}[theorem]{Conjecture}

\theoremstyle{definition}
\newtheorem{definition}[theorem]{Definition}
\newtheorem{example}[theorem]{Example}
\newtheorem{remark}[theorem]{Remark}

% === Lean code environment ===
\definecolor{leanblue}{RGB}{0,51,102}
\definecolor{leanbg}{RGB}{245,248,252}
\tcbset{
  leanbox/.style={colback=leanbg, colframe=leanblue!60,
    fonttitle=\bfseries\ttfamily, title=#1,
    boxrule=0.5pt, arc=2pt, left=4pt, right=4pt}
}

% === Macros ===
\newcommand{\BISH}{\textsf{BISH}}
\newcommand{\WKL}{\textsf{WKL}}
\newcommand{\DC}{\textsf{DC}}
\newcommand{\WLPO}{\textsf{WLPO}}
\newcommand{\LLPO}{\textsf{LLPO}}
\newcommand{\LPO}{\textsf{LPO}}
\newcommand{\MP}{\textsf{MP}}
\newcommand{\CLASS}{\textsf{CLASS}}
\newcommand{\CRM}{\textsf{CRM}}
\newcommand{\Z}{\mathbb{Z}}
\newcommand{\Q}{\mathbb{Q}}
\newcommand{\R}{\mathbb{R}}
\newcommand{\N}{\mathbb{N}}

\lstset{
  basicstyle=\ttfamily\small,
  keywordstyle=\color{leanblue}\bfseries,
  commentstyle=\color{gray},
  breaklines=true,
  frame=none,
  columns=flexible,
  keepspaces=true
}

% ============================================================
\begin{document}

\title{\textbf{The Diagnostic Penumbra:\\
Constructive Reverse Mathematics of Bayesian Clinical Decision-Making}\\[6pt]
\large Paper~104, Clinical Sub-Series Paper~B,\\
of the Constructive Reverse Mathematics Series}

\author{Paul Chun-Kit Lee\\
\textit{New York University, Brooklyn, New York}\\
\texttt{dr.paul.c.lee@gmail.com}}

\date{March 2026}
\maketitle

% ============================================================
\begin{abstract}
Every diagnostic test produces a finite contingency table --- true positives, false positives, false negatives, true negatives --- from which sensitivity, specificity, and predictive values are computed as rational numbers ($\BISH$).  The clinical decision invokes Bayes' theorem: the posterior probability $P(D \mid +)$ is compared against a treatment threshold~$\tau$ to determine whether to treat.  When prevalence is a rational cohort statistic, the posterior is rational and the decision is $\BISH$-decidable.  When prevalence is a real-valued population parameter, the posterior passes through the continuum, and comparison with~$\tau$ requires Markov's Principle ($\BISH + \MP$) --- structurally identical to the $p$-value comparison in Paper~103.  We identify the \emph{Diagnostic Penumbra}: the grey zone where the posterior lies within~$\delta$ of the treatment threshold, where $\delta$ combines prevalence uncertainty and analytical imprecision.  For ROC curve analysis, comparing AUCs of two tests requires $\LPO$ (integral over continuous ROC); discrete trapezoidal AUC is $\BISH$.  The treatment threshold from continuous utility functions requires $\WKL$ (bounded optimisation).

The diagnostic pipeline decomposes as $7\;\BISH + 1\;\BISH{+}\MP + 1\;\WKL + 1\;\LPO = 10$ stages (70\%~$\BISH$) --- isomorphic to Paper~103.  Applied to the ESC 0/1h high-sensitivity troponin algorithm: 24.7\% of chest pain presentations (1{,}972 of 7{,}998 in the Mueller validation cohort) fall in the grey zone (5--52~ng/L), with an event rate of 7.1\% that sits squarely in the treatment-threshold range for invasive investigation.  Applied to D-dimer for PE exclusion and PSA screening, the same structure recurs: rational data is $\BISH$; the grey zone emerges at the interface with the continuum.  All results are formalised in Lean~4 with Mathlib (10 module files, 0 sorries on main theorems, 4 documentary axioms).  This is the second paper in the CRM Clinical Sub-Series.
\end{abstract}

% ============================================================
\section{Clinical Premise}\label{sec:premise}

A cardiologist receives a 58-year-old man in the emergency department with acute chest pain.  The ECG is non-diagnostic.  The first high-sensitivity troponin result returns at 12~ng/L --- above the rule-out threshold of 5~ng/L but below the rule-in threshold of 52~ng/L.  The patient is in the \emph{observe zone}.

The clinician must decide: admit for serial troponins and possible angiography, or discharge with outpatient follow-up?  The decision depends on the posterior probability of acute coronary syndrome given the test result, compared against the treatment threshold determined by the relative costs of missed diagnosis versus unnecessary catheterisation.

Three questions arise:

\begin{enumerate}[nosep]
\item \textbf{What does ``in the grey zone'' actually mean?}  The contingency table data and the test characteristics are rational numbers, computable from finite counts.  But the population prevalence of ACS in chest pain presentations is a real-valued parameter estimated from epidemiological studies.  The posterior probability inherits this real-valuedness.  Comparing a real-valued posterior with a rational threshold is not a finite computation.

\item \textbf{Is the grey zone a statistical artefact or a logical structure?}  The ESC 0/1h algorithm's observe zone (5--52~ng/L) captures 24.7\% of patients.  We show this zone is not merely a practical inconvenience but a \emph{logical structure} --- the Diagnostic Penumbra --- where the decision is classically resolvable but not constructively certifiable.

\item \textbf{When does the grey zone collapse?}  With rational cohort prevalence (a specific study population), the posterior is rational and the decision is decidable in $\BISH$.  The grey zone is a property of the passage from finite data to the continuum, not of the data itself.
\end{enumerate}

These questions are the diagnostic analogues of the questions Paper~103 posed for the RCT statistical pipeline.  There, the transcendental gate was the normal CDF~$\Phi$; here, it is Bayes' theorem with real prevalence.  The mechanism is the same: rational data passes through a real-valued function, producing a quantity whose comparison with a rational threshold requires omniscience.

% ============================================================
\section{Introduction}\label{sec:intro}

\subsection{Main Results}

\begin{theorem}[A: Diagnostic Penumbra]\label{thm:A}
A diagnostic result with posterior $P(D \mid +)$ and total diagnostic uncertainty~$\delta = \delta_{\mathrm{prev}} + \delta_{\mathrm{cv}}$ lies in the \emph{Diagnostic Penumbra} at treatment threshold~$\tau$ if and only if
\[
  |P(D \mid +) - \tau| \;<\; \delta.
\]
The penumbra has width $2\delta$, symmetric around~$\tau$.  Within this band, the treat/no-treat decision is classically resolvable but not constructively certifiable.
\end{theorem}

\begin{theorem}[B: Bayesian MP Requirement]\label{thm:B}
When the prevalence~$\pi$ is a real-valued population parameter, the posterior
\[
  P(D \mid +) = \frac{\mathrm{sens} \cdot \pi}{\mathrm{sens} \cdot \pi + (1 - \mathrm{spec})(1 - \pi)}
\]
is a computable real.  Comparing $P(D \mid +)$ with a rational threshold~$\tau$ requires $\BISH + \MP$.  This is structurally identical to Paper~103's Theorem~B: rational data $\to$ transcendental function $\to$ real comparison $\to$ Markov's Principle.
\end{theorem}

\begin{theorem}[C: ROC = LPO]\label{thm:C}
Comparing AUCs of two diagnostic tests, where each AUC is the Riemann integral $\int_0^1 \mathrm{ROC}(t)\,dt$ over the continuous false-positive-rate axis, requires $\LPO$ (universal real comparison).  With discrete trapezoidal AUC from finitely many threshold points, the comparison is rational and $\BISH$-decidable.
\end{theorem}

\begin{theorem}[D: Discrete Bypass]\label{thm:D}
When the prevalence comes from a finite cohort study (rational), the entire Bayesian chain is $\BISH$:
\[
  \mathrm{sens}, \mathrm{spec}, \pi_{\mathrm{cohort}} \in \Q \;\Longrightarrow\; P(D \mid +) \in \Q \;\Longrightarrow\; \text{``$P(D \mid +) > \tau$?'' decidable in } \BISH.
\]
This is the diagnostic analogue of Fisher's exact test in Paper~103: stay in~$\Q$, pay no logical cost.
\end{theorem}

\subsection{CRM Primer}

Constructive Reverse Mathematics classifies mathematical theorems by the weakest axiom system above $\BISH$ (Bishop's constructive mathematics) required to prove them.  The hierarchy
\[
  \BISH \;\subset\; \BISH{+}\MP \;\subset\; \BISH{+}\LLPO \;\subset\; \BISH{+}\WLPO \;\subset\; \BISH{+}\LPO \;\subset\; \CLASS
\]
measures \emph{logical cost}: how much omniscience is needed.  Papers~1--102 established that the logical cost of mathematics is the logical cost of~$\R$ --- the completeness and total ordering at the Archimedean place~\cite{lee-crm-p98}.  Paper~103~\cite{lee-crm-p103} opened the Clinical Sub-Series by applying CRM to the RCT statistical pipeline.  Paper~104 extends the framework to diagnostic testing.

\subsection{Atlas Fit}

Within the CRM atlas, Paper~104 demonstrates the same \emph{Archimedean obstruction} as Paper~103 in a structurally distinct clinical domain:
\begin{itemize}[nosep]
\item \textbf{P103:} data $\xrightarrow{\text{CLT}}$ $\Phi$ $\xrightarrow{\text{compare}}$ $p < \alpha$ --- the transcendental gate is the normal CDF.
\item \textbf{P104:} data $\xrightarrow{\text{Bayes}}$ $P(D\mid+)$ $\xrightarrow{\text{compare}}$ $P > \tau$ --- the transcendental gate is Bayesian inference with real prevalence.
\end{itemize}
Both produce the same CRM decomposition: 7 BISH + 1 BISH+MP + 1 WKL + 1 LPO = 10 stages (70\% BISH).  The 70\% coincidence reflects the universal structure of the obstruction.

% ============================================================
\section{Preliminaries}\label{sec:prelim}

\subsection{Omniscience Principles}

We follow Bishop--Bridges~\cite{bishop-bridges} and Troelstra--van Dalen~\cite{troelstra-vandalen}.  A \emph{binary sequence} is a function $\alpha : \N \to \{0,1\}$.

\begin{definition}[LPO]
The \emph{Limited Principle of Omniscience}: for every binary sequence~$\alpha$, either $\exists\, n\,(\alpha(n)=1)$ or $\forall\, n\,(\alpha(n)=0)$.
\end{definition}

\begin{definition}[MP]
\emph{Markov's Principle}: if $\lnot\forall\, n\,(\alpha(n)=0)$, then $\exists\, n\,(\alpha(n)=1)$.
\end{definition}

\begin{definition}[WLPO]
The \emph{Weak Limited Principle of Omniscience}: for every~$\alpha$, either $\forall\, n\,(\alpha(n)=0)$ or $\lnot\forall\, n\,(\alpha(n)=0)$.
\end{definition}

\subsection{Decidability Hierarchy}

\begin{itemize}[nosep]
\item \textbf{Level~0 ($\BISH$):} Comparison of rationals $p, q \in \Q$ is decidable.
\item \textbf{Level~1 ($\BISH + \MP$):} Comparison of a computable real with a rational, given a $\lnot\lnot$-witness, terminates by Markov's Principle.
\item \textbf{Level~2 ($\LPO$):} Universal trichotomy on $\R$ is equivalent to LPO.
\end{itemize}

\subsection{Bayes' Theorem}

For a diagnostic test with sensitivity $\mathrm{sens}$ and specificity $\mathrm{spec}$ applied to a population with disease prevalence~$\pi$:
\begin{equation}\label{eq:bayes}
  P(D \mid +) = \frac{\mathrm{sens} \cdot \pi}{\mathrm{sens} \cdot \pi + (1 - \mathrm{spec})(1 - \pi)}.
\end{equation}
When $\mathrm{sens}, \mathrm{spec}, \pi \in \Q$, the posterior $P(D\mid +) \in \Q$.  When $\pi \in \R \setminus \Q$, the posterior inherits real-valuedness.

\subsection{Treatment Threshold}

The Pauker--Kassirer threshold~\cite{pauker-kassirer1980} for a treat/no-treat decision:
\begin{equation}\label{eq:threshold}
  \tau = \frac{C_{\mathrm{Rx}}}{B_{\mathrm{Rx}} + C_{\mathrm{Rx}}}
\end{equation}
where $B_{\mathrm{Rx}}$ is the net benefit of treating true disease and $C_{\mathrm{Rx}}$ is the net cost of treating non-disease.  With rational utilities, $\tau \in \Q$.

% ============================================================
\section{Main Results}\label{sec:main}

\subsection{The Contingency Table Is BISH}

\begin{proposition}\label{prop:ct-bish}
Every $2 \times 2$ contingency table $(TP, FP, FN, TN) \in \N^4$ produces rational test characteristics:
\[
  \mathrm{sens} = \frac{TP}{TP + FN} \in \Q, \quad
  \mathrm{spec} = \frac{TN}{TN + FP} \in \Q, \quad
  PPV = \frac{TP}{TP + FP} \in \Q, \quad
  NPV = \frac{TN}{TN + FN} \in \Q.
\]
All comparisons of these quantities (``Is sensitivity $> 0.90$?'') are decidable in $\BISH$.
\end{proposition}

\begin{proof}
Rational arithmetic is closed under division (with nonzero denominators).  Comparison of rationals is decidable by cross-multiplication.
\end{proof}

\subsection{The Bayesian Update: Rational vs Real}

\begin{proof}[Proof of Theorem~B]
The posterior~\eqref{eq:bayes} with $\mathrm{sens}, \mathrm{spec} \in \Q$ and $\pi \in \R$ is a composition of rational operations with one real input.  Since the function $f(\pi) = \mathrm{sens} \cdot \pi / [\mathrm{sens} \cdot \pi + (1-\mathrm{spec})(1-\pi)]$ is a rational function of~$\pi$, the output is a computable real when $\pi$ is computable.  By Bishop--Bridges~\cite{bishop-bridges}, \S2.3, comparing a computable real with a rational requires either an explicit separation witness (constructive) or Markov's Principle (to guarantee termination of the approximation).

The mechanism is identical to Paper~103's Theorem~B: there, $\Phi(z)$ was transcendental for algebraic~$z$ (Siegel--Shidlovskii); here, $f(\pi)$ is real-valued for real~$\pi$.  Both require $\BISH + \MP$ for comparison with a rational threshold.
\end{proof}

\begin{proof}[Proof of Theorem~D]
When $\pi = \pi_{\mathrm{cohort}} \in \Q$ (e.g., disease prevalence observed in a finite cohort study), the posterior $P(D\mid+) \in \Q$ by closure of $\Q$ under field operations.  Comparison with $\tau \in \Q$ is decidable by rational comparison.  No omniscience principle is invoked.
\end{proof}

\subsection{The Diagnostic Penumbra}

\begin{proof}[Proof of Theorem~A]
Define the total diagnostic uncertainty $\delta = \delta_{\mathrm{prev}} + \delta_{\mathrm{cv}}$, where $\delta_{\mathrm{prev}}$ is the half-width of the 95\% confidence interval on the prevalence estimate and $\delta_{\mathrm{cv}}$ is the analytical coefficient of variation of the assay scaled to the posterior.

The posterior $P(D\mid+)$ lies in the penumbra at threshold~$\tau$ when $|P(D\mid+) - \tau| < \delta$.  In this band, the classical value of $P(D\mid+)$ may be above or below~$\tau$, but the uncertainty~$\delta$ prevents constructive certification of which side.

The penumbra has width $2\delta$, symmetric around~$\tau$.  This is the diagnostic analogue of the Asymptotic Penumbra (P103), where width = $2\varepsilon$ with $\varepsilon$ the Studentized Berry--Esseen error.
\end{proof}

\subsection{ROC Analysis: Continuous vs Discrete}

\begin{proof}[Proof of Theorem~C]
The area under the ROC curve is $\mathrm{AUC} = \int_0^1 \mathrm{ROC}(t)\,dt$, a Riemann integral of a monotone nondecreasing function $[0,1] \to [0,1]$.  Comparing $\mathrm{AUC}_1 > \mathrm{AUC}_2$ for two diagnostic tests is a comparison of two real numbers.  By Bishop--Bridges~\cite{bishop-bridges}, universal real comparison $\forall x, y \in \R\,.\; x < y \lor x = y \lor x > y$ is equivalent to LPO.

The discrete bypass: when the ROC curve is approximated by finitely many operating points (each from a contingency table at a specific threshold), the trapezoidal AUC is a rational number, and comparison is decidable in $\BISH$.

This parallels the exact-vs-asymptotic dichotomy in P103: Fisher's exact test ($\BISH$, zero penumbra) vs asymptotic $z$-test ($\BISH+\MP$, nonzero penumbra).
\end{proof}

\subsection{Treatment Threshold: Rational vs Continuous}

When utilities are discrete and rational, the Pauker--Kassirer threshold~\eqref{eq:threshold} is rational and the decision is $\BISH$.  When utilities are defined over continuous outcome distributions (quality-adjusted life years, survival curves), the threshold becomes the root of an expected-utility equation over~$\R$, requiring at least $\WKL$ for bounded optimisation.

% ============================================================
\section{Clinical Test Cases}\label{sec:clinical}

\subsection{High-Sensitivity Troponin: ESC 0/1h Algorithm}

The ESC 0/1h algorithm~\cite{collet-esc2020,mueller2019} for acute chest pain uses high-sensitivity cardiac troponin I (hs-cTnI, Abbott Architect) with thresholds:
\begin{itemize}[nosep]
\item \textbf{Rule-out:} hs-cTnI $< 5$~ng/L at presentation $\Rightarrow$ discharge (NPV $> 99\%$)
\item \textbf{Rule-in:} hs-cTnI $\geq 52$~ng/L $\Rightarrow$ admit for invasive strategy (PPV $\approx 34\%$)
\item \textbf{Observe:} 5--52~ng/L $\Rightarrow$ serial troponin at 1~hour, clinical reassessment
\end{itemize}

\paragraph{Mueller validation cohort.}  In the international validation study~\cite{mueller2019} ($N = 7{,}998$):
\begin{itemize}[nosep]
\item Rule-out: $n = 4{,}593$ (57.4\%), 30-day MACE: 14 (0.3\%)
\item Observe: $n = 1{,}972$ (24.7\%), 30-day MACE: 141 (7.1\%)
\item Rule-in: $n = 1{,}433$ (17.9\%), 30-day MACE: 484 (33.8\%)
\end{itemize}

The observe zone captures \textbf{24.7\% of all chest pain presentations} --- nearly 1 in 4 patients.  The event rate of 7.1\% sits within the typical treatment-threshold range for invasive investigation ($\tau \approx 5$--15\% depending on utilities).  This is the Diagnostic Penumbra: the posterior is too close to the threshold for constructive resolution.

\paragraph{CRM analysis.}  The contingency table at the 52~ng/L threshold gives:
\begin{center}
\begin{tabular}{lcc}
\toprule
& MACE & No MACE \\
\midrule
hs-cTnI $\geq 52$ & 484 (TP) & 949 (FP) \\
hs-cTnI $< 52$ & 155 (FN) & 6{,}410 (TN) \\
\bottomrule
\end{tabular}
\end{center}

All derived quantities are rational: sensitivity $= 484/639 \approx 75.7\%$, specificity $= 6{,}410/7{,}359 \approx 87.1\%$, PPV $= 484/1{,}433 \approx 33.8\%$.  These computations are $\BISH$.

With the cohort prevalence $\pi = 639/7{,}998 \approx 8.0\%$, the Bayesian posterior for a positive test is $P(D\mid+) = 484/1{,}433 \approx 33.8\%$ (rational, $\BISH$-decidable).  With a population prevalence estimated from meta-analysis ($\pi \in \R$), the posterior inherits real-valuedness and comparison with~$\tau$ requires $\BISH + \MP$.

\subsection{D-dimer for Pulmonary Embolism}

The D-dimer assay at the standard threshold of 500~ng/mL (FEU) has sensitivity~$\approx 95\%$ and specificity~$\approx 40\%$ for PE~\cite{dinisio2007}.  At low pretest probability ($\pi \approx 5\%$, Wells $\leq 4$), the posterior probability of PE given a negative D-dimer is $\approx 0.65\%$ --- well below any treatment threshold.  The decision to exclude PE is $\BISH$-decidable even with real prevalence, because the separation from the threshold is so large that any rational approximation suffices.

Age-adjusted thresholds (age $\times 10$~ng/mL for patients $> 50$ years) reduce unnecessary imaging by approximately 30\%~\cite{righini2014} while maintaining the NPV above 97\%.

The D-dimer case demonstrates a \emph{collapsed penumbra}: the test's high sensitivity creates such strong posterior separation from the treatment threshold that the grey zone effectively vanishes for the rule-out application.  The penumbra persists for intermediate pretest probabilities.

\subsection{PSA Screening: The Degenerate Penumbra}

Prostate-specific antigen (PSA) at the traditional threshold of 4.0~ng/mL has sensitivity~$\approx 27\%$ for all-grade prostate cancer and specificity~$\approx 85\%$~\cite{thompson2004}.  The Thompson PCPT data revealed that 15.2\% of men with PSA $< 4.0$ harbour cancer.  With the standard threshold:
\[
  P(\text{cancer} \mid \text{PSA} \geq 4) \approx 31\%, \qquad P(\text{cancer} \mid \text{PSA} < 4) \approx 17.7\%.
\]
Both posteriors fall within the clinical grey zone for biopsy decisions (5--40\%).  The entire PSA range is a grey zone --- a \emph{degenerate penumbra} where no single threshold produces clean constructive separation.  This pathology motivated the shift from PSA-based decisions to risk calculators incorporating multiple variables.

% ============================================================
\section{CRM Audit}\label{sec:crm}

\subsection{Pipeline Classification}

\begin{table}[ht]
\centering
\caption{CRM pipeline for diagnostic testing}
\label{tab:pipeline}
\begin{tabular}{clc}
\toprule
Stage & Component & CRM Level \\
\midrule
1 & Contingency table (count data) & $\BISH$ \\
2 & Test characteristics (sens, spec, PPV, NPV) & $\BISH$ \\
3 & Likelihood ratios & $\BISH$ \\
4 & Bayesian update (rational prevalence) & $\BISH$ \\
5 & Bayesian update (real prevalence) & $\BISH + \MP$ \\
6 & Treatment threshold (rational utilities) & $\BISH$ \\
7 & Treatment threshold (continuous utilities) & $\WKL$ \\
8 & Clinical decision (rational posterior \& $\tau$) & $\BISH$ \\
9 & ROC AUC comparison (continuous) & $\LPO$ \\
10 & Test comparison (discrete ROC) & $\BISH$ \\
\midrule
\multicolumn{2}{l}{\textbf{Total: 7 BISH + 1 BISH+MP + 1 WKL + 1 LPO}} & \textbf{70\% BISH} \\
\bottomrule
\end{tabular}
\end{table}

\subsection{Comparison with Paper~103}

\begin{table}[ht]
\centering
\caption{Pipeline isomorphism between P103 and P104}
\label{tab:comparison}
\begin{tabular}{llll}
\toprule
Component & P103 mechanism & P104 mechanism & CRM \\
\midrule
Raw data & Sample statistics & Contingency table & $\BISH$ \\
Transcendental gate & Normal CDF $\Phi$ & Bayes($\pi_{\mathrm{real}}$) & $\BISH + \MP$ \\
Optimisation & Cox MLE & Utility threshold & $\WKL$ \\
Universal comparison & Real $p$-value & AUC integral & $\LPO$ \\
BISH fraction & 70\% & 70\% & --- \\
\bottomrule
\end{tabular}
\end{table}

The pipeline decompositions are isomorphic: both have 7 BISH + 1 BISH+MP + 1 WKL + 1 LPO.  The 70\% BISH rate is not a coincidence but reflects the universal structure of the Archimedean obstruction: most clinical computation is rational arithmetic over finite data; the non-constructive cost enters only at the interface between rational data and real-valued functions.

\subsection{What Descends from Where to Where}

\begin{itemize}[nosep]
\item $\BISH + \MP \to \BISH$: Bayesian inference with real prevalence descends to $\BISH$ when prevalence is rational (cohort study). \emph{Excision:} replace population prevalence with finite cohort data.
\item $\WKL \to \BISH$: Continuous utility threshold descends to $\BISH$ when utilities are discrete. \emph{Excision:} use the Pauker--Kassirer rational threshold model.
\item $\LPO \to \BISH$: Continuous ROC AUC descends to $\BISH$ with discrete operating points. \emph{Excision:} trapezoidal rule at finitely many thresholds.
\end{itemize}

Every non-$\BISH$ stage admits a $\BISH$ descent via discretisation.  The logical cost vanishes when the computation stays within~$\Q$.

% ============================================================
\section{Formal Verification}\label{sec:lean}

\subsection{Lean 4 Bundle Structure}

The formal verification comprises 10 module files:

\begin{tcolorbox}[leanbox={P104\_DiagnosticTesting/ — Module Structure}]
\begin{lstlisting}
OmnisciencePrinciples.lean    -- LPO, WLPO, MP definitions and implications
ContingencyTable.lean         -- 2x2 table, sens, spec, PPV, NPV
BayesianInference.lean        -- Rational/real posterior, MP requirement
ROCAnalysis.lean              -- AUC as integral, discrete bypass
TreatmentThreshold.lean       -- Pauker-Kassirer, rational/continuous
DiagnosticPenumbra.lean       -- Grey zone characterisation
TroponinCalculations.lean     -- ESC 0/1h algorithm verification
ClinicalTestCases.lean        -- D-dimer, PSA test cases
CRMPipeline.lean              -- Pipeline classification table
Paper104_Assembly.lean         -- Master theorem, axiom documentation
\end{lstlisting}
\end{tcolorbox}

\subsection{Key Code Snippets}

\paragraph{Bayesian posterior (rational).}

\begin{tcolorbox}[leanbox={BayesianInference.lean — Discrete Bypass}]
\begin{lstlisting}
theorem discrete_bypass (sens spec pi_cohort tau : Q)
    (h_pi_pos : 0 < pi_cohort) ... :
    Decidable (bayesPosteriorRat sens spec pi_cohort ... > tau) :=
  inferInstance
\end{lstlisting}
\end{tcolorbox}

The proof is \texttt{inferInstance}: Lean's type class inference finds the \texttt{DecidableEq} instance for~$\Q$ automatically.  No axioms invoked.

\paragraph{Pipeline classification.}

\begin{tcolorbox}[leanbox={CRMPipeline.lean — Pipeline Statistics}]
\begin{lstlisting}
theorem bish_count : bishStages.length = 7 := by native_decide
theorem non_bish_count : nonBishStages.length = 3 := by native_decide
theorem bish_percentage : 100 * bishStages.length / allStages.length = 70
  := by native_decide
\end{lstlisting}
\end{tcolorbox}

\paragraph{Troponin verification.}

\begin{tcolorbox}[leanbox={TroponinCalculations.lean — Grey Zone Fraction}]
\begin{lstlisting}
theorem observe_zone_approx :
    (1972 : Q) / 7998 > 24 / 100 /\ (1972 : Q) / 7998 < 25 / 100 := by
  constructor <;> norm_num
\end{lstlisting}
\end{tcolorbox}

\subsection{Axiom Inventory}

\begin{table}[ht]
\centering
\caption{Documentary axioms in Paper~104}
\label{tab:axioms}
\begin{tabular}{lll}
\toprule
Axiom & Purpose & Reference \\
\midrule
\texttt{real\_posterior\_comparison\_requires\_MP} & Posterior comparison needs MP & Bishop--Bridges \S2.3 \\
\texttt{auc\_comparison\_requires\_LPO} & AUC comparison needs LPO & Bishop--Bridges \S2.3 \\
\texttt{continuous\_threshold\_exists} & Continuous $\tau^*$ existence & Pauker--Kassirer \\
\texttt{continuous\_threshold\_requires\_WKL} & Continuous $\tau^*$ needs WKL & Paper~23 \\
\bottomrule
\end{tabular}
\end{table}

\subsection{Classical.choice Audit}

All theorems over~$\R$ show \texttt{Classical.choice} in \texttt{\#print axioms} --- this is Lean infrastructure (required for real number arithmetic in Mathlib), not classical mathematical content.  The constructive stratification is established by proof content: the main theorem \texttt{paper104\_main} uses only \texttt{inferInstance} (decidable instances), \texttt{norm\_num} (rational arithmetic), and \texttt{native\_decide} (finite enumeration).  No classical choice principle is invoked as mathematical content.

% ============================================================
\section{Discussion}\label{sec:discussion}

\subsection{The Archimedean Obstruction in Medicine}

Paper~103 identified the Archimedean obstruction in the RCT pipeline: rational test statistics pass through the transcendental normal CDF, producing real $p$-values whose comparison with rational thresholds requires $\MP$.  Paper~104 identifies the same obstruction in diagnostic testing: rational test characteristics pass through Bayes' theorem with real prevalence, producing real posteriors whose comparison with rational treatment thresholds requires $\MP$.

The mechanism is identical in both cases:
\[
  \text{rational data} \;\xrightarrow{\text{transcendental function}}\; \text{real output} \;\xrightarrow{\text{compare with rational}}\; \text{decision}.
\]
The first arrow is $\BISH$; the comparison requires $\MP$.  This is the statistical manifestation of the Archimedean Obstruction Thesis (Paper~98): the non-constructive cost enters through the interface between discrete data and the real continuum.

\subsection{Clinical Implications}

The Diagnostic Penumbra provides a new quality axis for evaluating diagnostic tests, orthogonal to sensitivity, specificity, and predictive values:

\begin{enumerate}[nosep]
\item \textbf{Troponin (ESC 0/1h):} Penumbra captures 24.7\% of patients.  The observe zone is \emph{logically necessary}, not merely a conservative clinical choice.
\item \textbf{D-dimer:} Penumbra collapses for the rule-out application at low pretest probability.  High sensitivity creates sufficient posterior separation.
\item \textbf{PSA:} Degenerate penumbra --- the entire range is a grey zone.  Poor sensitivity destroys posterior separation.
\end{enumerate}

\subsection{The Discrete Bypass Principle}

Every non-$\BISH$ stage in the diagnostic pipeline admits a $\BISH$ descent via discretisation (Theorem~D).  This is not an approximation; it is an exact logical downgrade.  When the clinician uses a cohort prevalence (rational) rather than a population prevalence (real), the posterior is exactly rational and the decision is exactly $\BISH$-decidable.

The clinical implication: \emph{institutional data is logically cheaper than epidemiological estimates}.  A hospital's own ACS prevalence from last year's chest pain presentations gives a rational prevalence that makes the Bayesian chain $\BISH$.  A meta-analytic estimate from 50 international studies gives a real-valued parameter that lifts the chain to $\BISH + \MP$.

% ============================================================
\section{Conclusion}\label{sec:conclusion}

The diagnostic testing pipeline decomposes into 10 stages: $7\;\BISH + 1\;\BISH{+}\MP + 1\;\WKL + 1\;\LPO = 10$ stages (70\%~$\BISH$), isomorphic to the RCT pipeline of Paper~103.  The Diagnostic Penumbra --- the grey zone where the posterior is too close to the treatment threshold for constructive resolution --- is the structural twin of the Asymptotic Penumbra.  Both arise from the same mechanism: rational data passing through the continuum.

Applied to troponin, 24.7\% of chest pain presentations fall in the penumbra.  The ESC 0/1h observe zone is not a clinical convenience but a logical necessity.  Applied to D-dimer, the penumbra collapses for rule-out at low pretest probability.  Applied to PSA, the penumbra degenerates to cover the entire range, reflecting the test's fundamental inability to separate.

The Clinical Sub-Series has now demonstrated the Archimedean obstruction in two independent clinical domains: RCT statistics (P103) and diagnostic testing (P104).  The 70\%~$\BISH$ rate --- identical in both --- is the signature of the universal architecture: most of medicine is rational arithmetic; the logical cost is concentrated at the interface with the continuum.

% ============================================================
\section*{Acknowledgments}

This paper was written with AI assistance (Anthropic Claude) for Lean~4 code generation, LaTeX formatting, and editorial review.  All mathematical content, clinical analysis, and CRM classifications are the author's.

Portions of the Lean formalisation use Mathlib~\cite{mathlib}, the mathematical library for Lean~4.

This paper follows the format guide for the CRM series~\cite{lee-format-guide}.

% ============================================================
\begin{thebibliography}{34}

\bibitem{bishop-bridges}
E.~Bishop and D.~Bridges, \emph{Constructive Analysis}, Springer, 1985.

\bibitem{troelstra-vandalen}
A.~S.~Troelstra and D.~van Dalen, \emph{Constructivism in Mathematics: An Introduction}, North-Holland, 1988.

\bibitem{pauker-kassirer1980}
S.~G.~Pauker and J.~P.~Kassirer, ``The threshold approach to clinical decision making,'' \emph{N.~Engl.~J.~Med.} \textbf{302} (1980), 1109--1117.

\bibitem{fagan1975}
T.~J.~Fagan, ``Letter: Nomogram for Bayes theorem,'' \emph{N.~Engl.~J.~Med.} \textbf{293} (1975), 257.

\bibitem{pencina2008}
M.~J.~Pencina, R.~B.~D'Agostino, R.~B.~D'Agostino~Jr, and R.~S.~Vasan, ``Evaluating the added predictive ability of a new marker: from area under the ROC curve to reclassification and beyond,'' \emph{Stat.~Med.} \textbf{27} (2008), 157--172.

\bibitem{sox2013}
H.~C.~Sox, M.~C.~Higgins, and D.~K.~Owens, \emph{Medical Decision Making}, 2nd ed., Wiley, 2013.

\bibitem{apple2021}
F.~S.~Apple et~al., ``Cardiac troponin assays: guide to understanding analytical characteristics and their impact on clinical care,'' \emph{Clin.~Chem.} \textbf{67} (2021), 136--148.

\bibitem{collet-esc2020}
J.-P.~Collet et~al., ``2020 ESC Guidelines for the management of acute coronary syndromes in patients presenting without persistent ST-segment elevation,'' \emph{Eur.~Heart~J.} \textbf{42} (2021), 1289--1367.

\bibitem{mueller2019}
C.~Mueller et~al., ``Multicenter evaluation of a 0-hour/1-hour algorithm in the diagnosis of myocardial infarction with high-sensitivity cardiac troponin~T,'' \emph{Ann.~Emerg.~Med.} \textbf{72} (2018), 654--665.

\bibitem{reichlin2009}
T.~Reichlin et~al., ``Early diagnosis of myocardial infarction with sensitive cardiac troponin assays,'' \emph{N.~Engl.~J.~Med.} \textbf{361} (2009), 858--867.

\bibitem{dinisio2007}
M.~Di~Nisio et~al., ``Diagnostic accuracy of D-dimer test for exclusion of venous thromboembolism: a systematic review,'' \emph{J.~Thromb.~Haemost.} \textbf{5} (2007), 296--304.

\bibitem{righini2014}
M.~Righini et~al., ``Age-adjusted D-dimer cutoff levels to rule out pulmonary embolism: the ADJUST-PE study,'' \emph{JAMA} \textbf{311} (2014), 1117--1124.

\bibitem{thompson2004}
I.~M.~Thompson et~al., ``Prevalence of prostate cancer among men with a prostate-specific antigen level $\leq$ 4.0~ng per milliliter,'' \emph{N.~Engl.~J.~Med.} \textbf{350} (2004), 2239--2246.

\bibitem{mathlib}
The Mathlib Community, \emph{Mathlib4}, \url{https://github.com/leanprover-community/mathlib4}, 2024.

\bibitem{lee-format-guide}
P.~C.~Lee, ``CRM series format guide,'' Zenodo, 2025.  DOI: 10.5281/zenodo.18765700.

\bibitem{lee-crm-p98}
P.~C.~Lee, ``The Motivic CRM Architecture,'' Paper~98 of the CRM Series, Zenodo, 2026.  DOI: 10.5281/zenodo.18828345.

\bibitem{lee-crm-p103}
P.~C.~Lee, ``The Asymptotic Penumbra: CRM of the RCT Statistical Pipeline,'' Paper~103 of the CRM Series, Zenodo, 2026.  DOI: 10.5281/zenodo.18880102.

\bibitem{lee-crm-p50}
P.~C.~Lee, ``The CRM Atlas,'' Paper~50 of the CRM Series, Zenodo, 2025.

\bibitem{lee-crm-p102}
P.~C.~Lee, ``Conservation Theorems for CRM,'' Paper~102 of the CRM Series, Zenodo, 2026.

\bibitem{series}
P.~C.~Lee, \emph{Constructive Reverse Mathematics Series}, Papers 1--103, Zenodo, 2024--2026.  \url{https://github.com/AICardiologist/FoundationRelativity}.

\bibitem{bentkus2003}
V.~Bentkus, ``On the dependence of the Berry--Esseen bound on dimension,'' \emph{J.~Statist.~Plann.~Inference} \textbf{113} (2003), 385--402.

\bibitem{pepe2003}
M.~S.~Pepe, \emph{The Statistical Evaluation of Medical Tests for Classification and Prediction}, Oxford University Press, 2003.

\bibitem{zweig-campbell1993}
M.~H.~Zweig and G.~Campbell, ``Receiver-operating characteristic (ROC) plots: a fundamental evaluation tool in clinical medicine,'' \emph{Clin.~Chem.} \textbf{39} (1993), 561--577.

\bibitem{sackett2000}
D.~L.~Sackett et~al., \emph{Evidence-Based Medicine: How to Practice and Teach EBM}, 2nd ed., Churchill Livingstone, 2000.

\bibitem{altman-bland1994}
D.~G.~Altman and J.~M.~Bland, ``Diagnostic tests~1: sensitivity and specificity,'' \emph{BMJ} \textbf{308} (1994), 1552.

\bibitem{deeks-altman2004}
J.~J.~Deeks and D.~G.~Altman, ``Diagnostic tests~4: likelihood ratios,'' \emph{BMJ} \textbf{329} (2004), 168--169.

\bibitem{moons2012}
K.~G.~M.~Moons et~al., ``Risk prediction models: I.\ Development, internal validation, and assessing the incremental value of a new (bio)marker,'' \emph{Heart} \textbf{98} (2012), 683--690.

\bibitem{vickers2006}
A.~J.~Vickers and E.~B.~Elkin, ``Decision curve analysis: a novel method for evaluating prediction models,'' \emph{Med.~Decis.~Making} \textbf{26} (2006), 565--574.

\end{thebibliography}

\end{document}
